\subsection{NN-G032: bounded-mass lift existence is fixed-parameter tractable}
\paragraph{Statement.}
For the symmetric effective associative lift problem of NN-G029,
let $N$ be a bound on the total coefficient mass, including
the unit, and let $L$ be the bit length of the sparse input
trace table, its rank and its specified unit and forbidden
positions, and the budget. There are a computable function
$F$ and an absolute constant $c$ such that lift existence
with mass at most $N$ is decidable in time
\[
                    F(N)(L+1)^c.                  \tag{G032.1}
\]
Thus the problem is fixed-parameter tractable in the numerical
mass budget. This is valid even when the budget is encoded
in binary; it is not a polynomial-time claim for unbounded
$N$. A lift can also be produced in the same type of bound.
The claim concerns the unrestricted symmetric table problem;
it does not identify that optimization with a canonical GCT
basis or imply any positive multiplicity formula.

The proof eliminates integer exponent variables symbolically.
The new consequence is the parameterized algorithm for this
lift problem; the integer-linear quantifier elimination used
below is a standard method, proved here self-containedly.
It never searches through all integers up to the largest input
trace exponent. Its constants can be extremely large, so it
does not make the $N=11147$ example a practical computation.

\paragraph{Finite discrete patterns and integer formulas.}
First canonically combine repeated sparse monomials, and
check the prescribed symmetry, unit and input format. These
operations take polynomial time in $L$. If the absolute trace
mass exceeds $N$, reject. A unital rank-$r$ lift contains at
least $2r-1$ unit coefficients, so reject also if $2r-1>N$.
Consequently every nonrejected input has $r\le(N+1)/2$,
at most $N$ nonzero trace monomials, and trace coefficient
absolute values at most $N$.

The following discrete information consequently has only
a number of possibilities bounded by a computable function
of $N$: the rank; the occupied trace positions and their
integer coefficients; the forbidden positions; and the
coefficient-and-position patterns of all possible effective
lift tables of total mass at most $N$.
For a lift pattern, every noncentral symmetric monomial pair
costs at least two units of mass. Assign a distinct formal
positive integer variable to each occurrence of such a pair,
and retain the central coefficients as fixed integers.
There are at most $\lfloor N/2\rfloor$ exponent variables.
Within each position, require that separate pairs have
distinct exponents. Coefficients are selected as part of
the finite pattern, so there are no unknown coefficient
products in what follows.

Let $k_1,\ldots,k_h$ be the distinct positive exponents of
the input trace. They are free integer parameters; their
binary magnitudes are not included in the discrete pattern.
One has $h\le\lfloor N/2\rfloor$.
As in NN-G029, form zero, the signed individual exponent
variables and trace parameters, and all sums of two signed
individual forms. Enumerate their possible total order
and equality patterns. Each chosen pattern is expressed
by a conjunction of linear equalities and strict
inequalities with integer coefficients bounded in terms
of $N$. Since the variables are integers, a strict
inequality is an ordinary weak inequality with an added
constant one.

For a fixed collision pattern, domination and parity are
comparisons of fixed coefficients, and every associativity
coefficient is a fixed integer sum of fixed integer products.
Discard a pattern if these tests fail, or if it violates a
unit or forbidden-position condition. Every remaining
pattern yields a formula
\[
 \exists x_1,\ldots,x_m\in\mathbb Z\quad
       \Phi(x_1,\ldots,x_m;k_1,\ldots,k_h),          \tag{G032.2}
\]
where $m\le N/2$ and $\Phi$ is a finite conjunction of
integer linear inequalities, including $x_i\ge1$.
The numbers of patterns, atoms and variables, and all
coefficients and constants in these formulas, are bounded
by functions of $N$ alone. The input trace exponents
appear only as free parameters $k_j$.

Conversely, a satisfying integer assignment to any retained
formula gives an actual lift: its monomials have the selected
positive exponents; its constants dominate the trace with
the correct parity; and all coefficient identities for
associativity hold by the selected collision partition.
Thus the finite disjunction of (G032.2) is equivalent to
the desired bounded-mass existence question. This step
does not require an incumbent or the numerical cutoff
from NN-G029.

\paragraph{Elementary elimination of one integer variable.}
We prove the needed elimination rule completely. Consider
any finite Boolean combination of integer linear inequalities
and congruences in one quantified integer $y$, other integer
variables, and free parameters. All coefficients and moduli
are fixed integers. Equalities are pairs of inequalities.
Negating an inequality changes it to an inequality with
constant shift one. A negated congruence modulo a fixed
positive $m$ is a disjunction over its other $m-1$ residues.
After distributing conjunctions and disjunctions, it is
enough to eliminate $y$ from one conjunction of inequalities
and congruences. This may greatly increase size, but the
increase depends only on the formula, not on parameter values.

Let $B\ge1$ be the least common multiple of the absolute
values of all nonzero coefficients of $y$ in that conjunction.
If none occurs, set $B=1$. Substitute $z=By$ and add
$z\equiv0\pmod B$.
An atom $ay+t\ge0$, where $t$ does not involve $y$,
becomes, for $a\ne0$,
\[
 \operatorname{sgn}(a)z+(B/|a|)t\ge0.
\]
A congruence $ay+t\equiv0\pmod m$ becomes
\[
 \operatorname{sgn}(a)z+(B/|a|)t
                   \equiv0\pmod{mB/|a|}.
\]
These transformations are equivalent under $z=By$:
the inequality is multiplied by a positive integer, and
$m\mid u$ is equivalent to $cm\mid cu$ for every
positive integer $c$. Atoms with $a=0$ are unchanged.
Conversely $z\equiv0\pmod B$ recovers an integer $y$.

All remaining atoms involving $z$ are therefore lower bounds
$z\ge\ell_i$, upper bounds $z\le u_j$, or congruences
$z\equiv r_a\pmod{m_a}$, where the bounds and residues
are integer linear forms in the other variables and free
parameters. Let $M$ be the least common multiple of the
fixed positive moduli. Enumerate $a=0,\ldots,M-1$ and
impose the compatibility atoms
\[
                      a\equiv r_b\pmod{m_b}
                 \quad\hbox{for every }b.
\]
These atoms exactly say that the chosen residue class
$z\equiv a\pmod M$ satisfies every congruence.

If either the lower-bound list or the upper-bound list is
empty, that residue class has an integer satisfying every
remaining one-sided bound. Thus no further condition is
needed in that case. If both lists are nonempty, enumerate
an index $i$ for a largest lower bound and a number
$b\in\{0,\ldots,M-1\}$, and impose
\[
 \ell_i\ge\ell_{i'}\ (\hbox{all }i'),\qquad
 \ell_i+b\equiv a\pmod M,\qquad
 \ell_i+b\le u_j\ (\hbox{all }j).
\]
These conditions are necessary and sufficient: the smallest
integer in the chosen residue class at least $\ell_i$ is
exactly $\ell_i+b$ for its unique such $b$.
The resulting finite disjunction contains no $z$ or $y$
and consists only of integer linear inequalities and
congruences. This proves elimination of one integer
existential quantifier, with an explicit finite procedure.

\paragraph{Runtime and witness recovery.}
Apply this elimination successively to all the exponent
variables in every retained formula (G032.2). Its initial
size, coefficients, moduli and number of variables depend
only on $N$. At every step the resulting size, coefficients
and moduli still depend only on $N$, even though they may
grow rapidly. The result is a finite Boolean combination
of integer linear inequalities and fixed-modulus
congruences in the input exponents $k_j$ alone.
Formula construction and elimination therefore cost some
computable function of $N$, independently of the binary
magnitudes of the $k_j$.

Substitute the given binary exponents into this formula.
There are only $F_0(N)$ atoms and coefficients, each of
bit length at most $F_0(N)$ for a computable $F_0$.
Integer addition, multiplication and division by these
fixed coefficients and moduli can all be carried out by
elementary bit algorithms in time bounded by
$F_1(N)(L+1)^2$. The finite initial parsing and sorting
contribute another fixed polynomial in $L$.
This proves (G032.1), with an absolute polynomial exponent.

One may also retain the selected disjunct at every elimination
step. Recover variables in reverse order, selecting a true
disjunct after substituting the free parameters and the variables
already recovered at that stage. With two-sided
bounds use the recorded integer $z=\ell_i+b$; with only
lower bounds choose the least compatible residue representative
above their maximum; with only upper bounds choose the
largest representative below their minimum; and with no
bounds choose the selected residue itself. Then set $y=z/B$.
The compatibility atoms ensure divisibility by $B$.
Repeating this construction yields the integer exponents
of a lift. The number of steps and every arithmetic
coefficient are bounded by functions of $N$, so recovery
obeys the same type of parameterized time and bit-length
bound. In particular every recovered bit length is
$O_N(L+1)$, since each recovery step uses a bounded number
of linear operations and fixed-modulus divisions with
coefficients depending only on $N$.\hfill$\square$

\paragraph{Geometric translation and limitations.}
For a fixed coefficient pattern, the grading weights lie
in the cocharacter lattice of a finitely presented
diagonalizable group, with the associativity collisions
imposing integral linear relations. Domination and the
chosen order of weights specify finitely many lattice
chambers. Eliminating exponent variables projects these
integer lattice conditions; the resulting congruences
record finite-index lattice obstructions. Their character
duals are finite diagonalizable groups such as $\mu_m$.
The actual conclusion is that, for bounded total dimension
mass, admissibility of the prescribed trace weights is
decidable by finitely many such linear inequalities and
congruences, computed uniformly from the mass budget.

This is an arithmetic statement about cocharacter lattices,
not a claim that inequalities define Zariski-open subsets
or that an effective lift is a canonical geometric object.
The grading dictionary is SGA 3, I, 4.7.2--4.7.3, and
the actual composing kernel realization is NN-G010,
using archive cote 2, scan 17. The algorithm gives no
polynomial bound on the mass budget needed by varying
GCT inputs, no practical small value of $F$, and no
rank-four positivity, purity, PH0 or PH1 theorem.

